\section{Data}
\label{sec:data}
The data used in this project consists of 2D brain MRI slices for a multi-class brain tumor classification and segmentation task. 
It comprises grayscale (single-channel) medical imaging data, where each sample is a 2D slice labeled as one of four categories—glioma, 
meningioma, pituitary tumor, or no tumor. The primary source is the publicly available Brain Tumor MRI Dataset hosted on Kaggle, containing 
approximately 7,200 human brain MRI images distributed across the four classes, supplemented by the SciDB Brain Tumor Dataset, which provides
 the T1-weighted contrast-enhanced images. We did not collect new data, and the project relies entirely on these two publicly available sources.

 The dataset was divided into stratified train/validation/test splits (80/10/10) to preserve class distribution and mitigate imbalance. For segmentation, the mask-annotated data are split into 1753 training images, 219 validation images, and 220 test images. For classification, the stratified split contains 5618 training images, 702 validation images, and 703 test images.

We addressed class imbalance through weighted cross-entropy loss based on inverse class frequency for classification, and a combined BCE--Dice loss for segmentation to handle the severe foreground-background disparity typical of tumor regions. During training, we apply online data augmentation: random spatial and intensity transformations with joint image-mask processing for segmentation, and random rotation, horizontal flipping, and brightness/contrast jittering for classification. No additional heavy filtering was required. Preprocessing mainly consists of consistent resizing, mask alignment, and online augmentation, all implemented in PyTorch.
    